\begin{figure}[htbp]
\centering
\begin{tikzpicture}
\begin{axis}[
  width=0.82\textwidth, height=6.0cm,
  xmin=1, xmax=128, ymin=0, ymax=45, xlabel={Number of cores $P$}, ylabel={Speedup $S$},
  xtick={1,8,16,32,64,96,128}, grid=major, grid style={draftgray!20},
  legend pos=north west, legend style={font=\small},
  every axis plot/.append style={thick,domain=1:128,samples=128}]
\addplot[draftblue] {1/(0.02+0.98/x)}; \addlegendentry{$f_{\mathrm{seq}} = 2\%$}
\addplot[draftblue,dashed] {1/(0.05+0.95/x)}; \addlegendentry{$f_{\mathrm{seq}} = 5\%$}
\addplot[draftgray,dotted] {1/(0.10+0.90/x)}; \addlegendentry{$f_{\mathrm{seq}} = 10\%$}
\addplot[only marks,mark=*,draftblue,mark size=1.6pt] coordinates {(8,5.93) (16,9.14) (64,15.42)};
\end{axis}
\end{tikzpicture}
\caption[Amdahl's law, , for three sequential fractions, after Sarangi]{Amdahl's law, \cref{eq:amdahl}, for three sequential fractions, after Sarangi \cite[sec.~12.2.3, Fig.~12.4]{sarangi2025basic}. The marked points are the values worked out for $f_{\mathrm{seq}} = 5\%$.}
\label{fig:amdahl}
\end{figure}
